%!TeX root=../main.tex

\palavraschave{Tutor inteligente, Ensino de shogi, Ensino adaptativo, Feedback automático}
\keywords{Intelligent tutoring system, Shogi learning, Adaptive learning, Automatic feedback}

\resumo{
Este trabalho tem como objetivo desenvolver o Shogi-sensei, uma ferramenta para 
ensinar os fundamentos básicos de shogi para iniciantes. Para tal, foi criado um 
gerador de exercícios organizados em módulos, cada um abordando as principais mecânicas do jogo. 
Além disso, foi implementado um algoritmo que adapta a distribuição dos módulos nas listas de exercícios
de acordo com o desempenho do usuário em cada um deles. Ao final das listas, um modelo de linguagem de
grande escala gera explicações sobre as questões respondidas incorretamente.
O servidor do sistema foi desenvolvido em Python, enquanto a interface foi implementada utilizando React. 
Para avaliar a efetividade do Shogi-sensei, foi realizado um experimento com participantes com pouco ou nenhum 
conhecimento de shogi, durante o  qual foram coletadas métricas como taxa de acerto e variação da distribuição de probabilidades
de seleção atribuídas a cada módulo. Ao final, os participantes também responderam a um questionário 
para compartilhar suas impressões sobre o sistema. Os resultados sugerem que a ferramenta 
contribuiu positivamente para o aprendizado dos fundamentos básicos do jogo. Dessa forma, conclui-se que,
embora ainda se encontre em um estágio inicial de desenvolvimento, o projeto atingiu seu objetivo de oferecer 
uma ferramenta de apoio ao aprendizado de shogi para iniciantes.
}

\abstract{
This work aims to develop Shogi-sensei, a tool designed to teach the basic fundamentals of shogi to beginners. 
To this end, an exercise generator was developed, organizing exercises into modules, each covering a different 
key mechanic of the game. Additionally, an algorithm was implemented to adapt the distribution of modules across 
exercise lists based on the user's performance in each module. At the end of each list, a large language model 
generates explanations for incorrectly answered questions.
The system's server was developed using Python, while the user interface was implemented using React. To evaluate 
the effectiveness of Shogi-sensei, an experiment was conducted with participants who had little or no prior knowledge 
of shogi. During the experiment, metrics such as accuracy and changes in the selection probability distribution assigned 
to each module were collected. At the end, participants also completed a questionnaire to share their impressions of the 
system. The results suggest that the tool had a positive impact on learning the basic fundamentals of the game. Thus, it 
is concluded that, although still at an early stage of development, the project achieved its initial goal of providing a 
tool to support shogi learning among beginners.
}
